\section{Companions that Gaia already excludes: MOLUSC}
\label{sec:molusc}

\subsection*{Why}
Several false-positive scenarios of TRICERATOPS need an unresolved
\emph{bound companion} of the target: PTP and PEB (the companion dilutes
the dip), STP and SEB (the dip is on the companion). Without other
information TRICERATOPS uses a general prior: how often stars of this mass
have a companion, and with which masses and separations (Moe \& Di~Stefano
2017). For many nearby stars that prior is too pessimistic: Gaia has
observed them for years, and a bright companion at the right distance would
have shown up. MOLUSC (\emph{Multi-Observational Limits on Unseen Stellar
Companions}, Wood, Mann \& Kraus 2021, AJ~162, 128) turns those
non-detections into a list of the companions that are still allowed. It
costs no telescope time.

\subsection*{How MOLUSC works}
\begin{enumerate}
\item \textbf{Simulate.} It draws a large number of companions (typically
$10^6$) for the target: mass ratio $q$ uniform between 0 and 1, a log-normal
period distribution (Raghavan et al.\ 2010), and random eccentricity,
inclination, orientation and orbital phase. Each simulated star has
exactly \emph{one} companion.
\item \textbf{Predict what each companion would have done}, with the
stellar models of Baraffe et al.\ (2015) for its brightness:
\begin{itemize}
\item \emph{Gaia astrometry (RUWE).} A companion at roughly 0.1--1$''$
makes the photocentre wobble and raises the RUWE (renormalised unit weight
error) of the Gaia solution. MOLUSC predicts the RUWE for every simulated
companion and rejects those whose prediction is clearly larger than the
observed value.
\item \emph{Gaia imaging.} A companion bright enough and far enough from the
star (beyond about 1$''$) would be a separate Gaia source.
\item \emph{Optional:} a high-resolution contrast curve (speckle, AO) and
radial-velocity measurements. With RVs, close companions (a few au and
less) can be excluded as well.
\end{itemize}
\item \textbf{Keep the survivors.} The file \code{<prefix>_kept.csv} lists
the companions that none of the data would have seen.
\end{enumerate}
The rejections are probabilistic (a companion is rejected with the
probability that it would have been detected), so a run with $10^6$
companions gives smooth distributions.

\subsection*{How the pipeline uses it}
\begin{enumerate}
\item \textbf{Run MOLUSC.} A copy with a few small fixes (newer SciPy/NumPy, a VizieR fallback) is
in the folder \code{MOLUSC} (see \code{README_pipeline.txt}). For each target a
short script gives the coordinates and the mass, e.g.\ for L~279-94:
\begin{lstlisting}
bash MOLUSC/run_molusc_L27994.sh       # in WSL, about 10 min
# -> TIC100796276_molusc/L27994_N1000000_kept.csv
\end{lstlisting}
MOLUSC needs the Gaia archive; when it does not answer (for instance while
it is closed for the DR4 release), the same Gaia DR3 values are read from
VizieR.
\item \textbf{Give the file to the pipeline:} GUI tab \emph{Field},
``MOLUSC companions file'' (\code{--molusc-file}) and ``MOLUSC: companions
generated'' (\code{--molusc-n}, the number MOLUSC simulated).
\item \textbf{What happens inside.} TRICERATOPS(+) takes the mass ratios of
the MOLUSC companions for the first draws of the PTP, PEB, STP and SEB
scenarios, and treats the remaining draws as ``no companion''. Two details
matter:
\begin{itemize}
\item It keeps only companions whose orbit never comes closer than 10~au to
the star (periastron $a(1-e) > 10$~au). Planets are rare in such close
binaries (Moe \& Kratter 2021), so TRICERATOPS assumes they do not exist
there.
\item The prior probability of a companion then becomes (number of
companions in the file)/(number of TRICERATOPS draws, $10^6$). Used as is,
the FPP would depend on how many companions one happened to simulate in
MOLUSC, and a file with more than $10^6$ survivors would crash
TRICERATOPS. The pipeline therefore resamples the survivors so that this
probability equals the MOLUSC survival fraction times an assumed
multiplicity (\code{--molusc-multiplicity}, default 1: every star has a
companion, the conservative choice).
\end{itemize}
The log lists the numbers: companions generated, kept, kept beyond 10~au,
the resulting prior, and their mass-ratio distribution. The audit names
the companion prior that was used (\emph{Evaluation}).
\end{enumerate}

\subsection*{Limitations}
\begin{itemize}
\item \textbf{It removes only what Gaia could have seen.} Companions very
close to the star (periastron $\lesssim$1~au) hardly change the RUWE and are
not resolved; without RV data they all survive. That they do not count in
the FPP is the \emph{assumption} that planets avoid close binaries, not a
measurement. A result that relies on MOLUSC should say so, and spectroscopy
(RVs or a second set of lines) is the observation that tests it.
\item \textbf{The RUWE is not only a binary indicator.} Variability, a
nearby source or a poor fit also raise it; a high RUWE makes MOLUSC reject
\emph{fewer} companions, which is safe. A RUWE that is low by chance would
make it reject too many.
\item \textbf{Models and distributions.} The companion population
(Raghavan), the stellar models (BHAC15, default age 5~Gyr) and the RUWE
calibration of the MOLUSC tables are assumptions; for very young or
evolved stars they matter.
\item \textbf{One companion per star.} Triples are not simulated. The
multiplicity (fraction of stars that have a companion at all) is not part
of MOLUSC; the pipeline uses 1 unless told otherwise, which overestimates
the companion probability (about a quarter of M dwarfs have one,
Winters et al.\ 2019).
\item \textbf{Different versions, different answers.} TRICERATOPS 1.1.0
treats the MOLUSC companions differently from the TRICERATOPS+ version
used here; the Zenodo note on L~279-94 obtains FPPs between 0.003 and 0.042
depending on such choices. Report the FPP with the companion model.
\item \textbf{Not a replacement for high-resolution imaging.} MOLUSC shows
that the companions which would invalidate the planet are \emph{unlikely};
an image or RVs show whether one is \emph{there}.
\end{itemize}
